\documentclass[11pt]{ltjsarticle}
\usepackage{luatexja-fontspec}
\usepackage{amsmath,amssymb}
\usepackage[margin=25mm]{geometry}
\usepackage{xcolor}
\usepackage{hyperref}
\hypersetup{colorlinks=true,linkcolor=blue,urlcolor=blue,citecolor=blue}
\setcounter{secnumdepth}{0}
\setlength{\parindent}{0pt}
\setlength{\parskip}{0.5em}
\providecommand{\tightlist}{\setlength{\itemsep}{0pt}\setlength{\parskip}{0pt}}
\begin{document}
\section{高い対称性の極限におけるスケールの対数写像とスケール・量子化の分離について}\label{ux9ad8ux3044ux5bfeux79f0ux6027ux306eux6975ux9650ux306bux304aux3051ux308bux30b9ux30b1ux30fcux30ebux306eux5bfeux6570ux5199ux50cfux3068ux30b9ux30b1ux30fcux30ebux91cfux5b50ux5316ux306eux5206ux96e2ux306bux3064ux3044ux3066}

\subsection{------最小公理系からの限定的検証報告（改訂版）}\label{ux6700ux5c0fux516cux7406ux7cfbux304bux3089ux306eux9650ux5b9aux7684ux691cux8a3cux5831ux544aux6539ux8a02ux7248}

\textbf{著者：} 木原 範昭（ORCID: 0009-0004-6753-4020） \textbf{日付：}
2026年6月18日（改訂版） \textbf{DOI（本版）：} 10.5281/zenodo.20740842
\textbf{Concept DOI（全版）：} 10.5281/zenodo.20740841 \textbf{Zenodo:}
https://zenodo.org/records/20740842

\begin{center}\rule{0.5\linewidth}{0.5pt}\end{center}

\subsection{読者への注記（最初に読んでください）}\label{ux8aadux8005ux3078ux306eux6ce8ux8a18ux6700ux521dux306bux8aadux3093ux3067ux304fux3060ux3055ux3044}

本稿は、新たな一般的事実の「発見」を主張するものではない。本稿は限定的な前提から出発するため、一読で誤読されやすい。誤読のほとんどは、\textbf{系の中（不可知・無名の内部視点）と系の外（観測者の視点）を曖昧に混ぜて読むこと}から生じる。なお本系では「どちらが実在か」という問いは成立しない（§0.5
の注）。客観的な錨は実在ではなく、不動点と保存量に置く。これを防ぐため、本稿は各言明の頭に、それが系の中の言明か系の外の言明かを原則として明示する。先に
§0.5（系の中と外）と §1（規約）を読まれたい。

\begin{center}\rule{0.5\linewidth}{0.5pt}\end{center}

\subsection{要旨（Abstract）}\label{ux8981ux65e8abstract}

本報告は、共役な二量の関係
νλ=k（観測者が読む量）への素朴な疑問から出発する。この硬い関係だけからは構造が生まれない。構造が入りうる唯一の隙間として、整数値
ν
が内在の揺らぎを持つこと、および系が高い無名性と高い対称性を持つことを置く。極めて高い対称性と無名性を要求する限定された極限において、スケールを表す二乗量が対数変換による接続を通じて平坦な加法構造（平坦な超平面）に写像され得る、という\textbf{限定範囲での検証結果}を報告する。

基本量を二乗量の対数 N=log\textsubscript{2}ν として導入すると、共役関係 νλ=k
が加法的恒等式 N\_ν+N\_λ=K となり、スケールを表す二乗和 Σν²=R²（g
側）が、自己双対の不動点（k=1）かつ全対称の極限において、平坦な一乗和
ΣN=K
に写像される。重要なのは、この対数写像がスケール（二乗和・g）を平坦化して
gauge として消す一方で、揺らぎ
Δ（積・面積・ω）はその平坦面上に不変に残ることである。すなわち本写像は、スケールと量子化を分離する。スケールは
g に宿り平坦化可能、量子化は ω に宿り写像で不変である。ここで R
は系のスケール（大きさ）であって、曲率テンソルでも球面の半径でもない。

本稿は揺らぎの大きさ ±1/2 を\textbf{導出ではなく措定}として扱う（理由は
§6）。本結果が既知の枠組み（Weyl
変換、共形平坦性等）の特殊例に留まるか否かは確定できず、反証の網羅的検証は行えていない。本報告の目的は、この限定的構造の記録に留まる。

\begin{center}\rule{0.5\linewidth}{0.5pt}\end{center}

\subsection{§0
動機（著者ノート）}\label{ux52d5ux6a5fux8457ux8005ux30ceux30fcux30c8}

本稿の出発点は、共役な二量の関係 νλ=k への素朴な疑問である。

【系の外（観測者）】νλ=k と置くと λ=k/ν、すなわち ν が大きくなれば λ
が小さくなる。これは外から観測される、硬いシーソーの関係である。

【系の外（観測者）】しかしこの硬い関係だけでは、ν と λ
という二量から構造が生まれない。比例関係が一本あるにすぎない。

では、構造が入りうる余地はどこにあるか。本稿は二つの緩めどころを置く。

【離散構造】第一に、ν
は離散（整数値）であり、連続な小数部を持たない。「本当は連続値があり、観測精度の限界で整数に見える」のではなく、系が離散だから整数である（離散性の導出は
§3）。これは「ν
が実在する」という主張ではない------本系では何が実在かを問えない（§0.5
の注）。整数性は実在の主張ではなく、離散構造の帰結である。ただし整数値は確定しておらず、内在の揺らぎを持つ。これは観測による誤差ではなく、系に内在する揺らぎである。

【系の中】第二に、νλ=k において、ν と λ
を区別する名前が無い。どちらがどちらか、系の中からは区別できない。すなわち本系は高い無名性を持つ。

本稿はあえて、この高い無名性と高い対称性を前提に置いて系を組み立てる。通常、これほど強い無名性と対称性を課すと自由度が拘束され、構造は生まれにくくなる。それでもなお何が立ち上がるかを見るのが、本稿の狙いである。

【系の外（観測者）／系の中】最後に一点。νλ=k
は、あくまで系の外の観測者の目線で読まれる量である。これとは別に、系の中に内在する不変量（保存量）あるいは不変点（安定点）があるのではないか------これが本稿を貫く問いである。

\begin{center}\rule{0.5\linewidth}{0.5pt}\end{center}

\subsection{§0.5
著者ノート：系の中と外（不可知と不変量）}\label{ux8457ux8005ux30ceux30fcux30c8ux7cfbux306eux4e2dux3068ux5916ux4e0dux53efux77e5ux3068ux4e0dux5909ux91cf}

本稿は全体を通じて系の中と系の外を厳密に区別する。ここを曖昧にすると、以下のほとんどの言明が誤読される。

\begin{itemize}
\tightlist
\item
  【系の中】系の中とは、無名（名も色も位置も持たない）の要素の視点である。系の中で区別できないものは、系の中では
  null（無）と等価に扱う。系の中では、系の外の量を直接観測することは原則できない（ただし系外の量が系内に影響として現れる場合は、その影響を通じて認知できる。例は注1）。
\item
  【系の外】系の外とは、外から観測する者の視点である。観測者が得られるのは限られた観測量だけで、系の中の不変量（不動点・保存量）に直接届くとは限らない。ゆえに観測者は、系の中に何があるか（あるいは無いか）を断定できない。
\end{itemize}

\textbf{【重要な注：実在は問えない】}
本系は双対・共役な関係（νλ=k）の上に立つ。ゆえに「ν と λ
のどちらが実在か」という問いは成立しない------一方を実在とすると、他方の実在が曖昧になる。系の中の実在を否定もしない。意味を持つのは「どちらが実在か」ではなく、双対のもとで不変に保たれるもの------\textbf{不動点（k=1,
ν=1）と保存量（積 δ\textsubscript{1}δ\textsubscript{2}=1/2、二乗和
Σν²）}------である。以下、客観的な錨はすべて不動点・保存量として述べ、実在の所在は主張しない。

この区別の下で、最も基本的な計数を系の中の視点から導く。系の中で登録できるのは「自分以外があるか、ないか」ただ一つである。

\begin{itemize}
\tightlist
\item
  【系の中】0 の場合：状態を判断する主体が無い。ゆえに 0 は null
  と区別できない。これは系の中の区別不能を述べるだけで、整数構造としての
  0 を排除しない。
\item
  【系の中】1
  の場合：唯一の要素には「自分以外」が無い。無名ゆえ、他者との対照なしには「自分が在る」ことさえ立たない。ゆえに系の中では
  1 は 0 と区別できない。これは整数構造としての 1 を排除しない。
\item
  【系の中】2
  の場合：初めて「自分以外」が在る。それを通じて「自分が在る」が立つ。2
  が、系の中から何かが知れる最小である。
\item
  【系の中】2
  以上の場合：他者は色も名も位置も持たないので数えられない。ゆえに系の中では
  2 と 2 以上は区別できない。これは整数構造としての 2 以上を排除しない。
\end{itemize}

【系の中】したがって系の中の視点は、二値------「他者なし」（0 と 1
を区別できない）と「他者あり」（2 と 2
以上を区別できない）------に畳まれる。

【保存量】一方、保存量・不動点の構造としては整数列 0,1,2,3,\ldots{}
が立つ。系の中の二値への畳み込みは、不可知による畳み込みであって、この整数構造を否定しない。

【系の中／外】要点を一行で：系の中で不可知なものは、系の中では null
と等価である。だがそれを系の外において否定もできない------不可知なものは、否定の対象にすらならない。

\textbf{注1（系外の量が系内に影響として現れる例）：}
多次元時空における空間曲率がその例である。曲率そのものは系の外の量で、系の中ではその方向を認知できないが、その影響------例えば三角形の内角の和が
180
度からずれること------を通じて系の中で認知できる。ただし本稿の系ではこの例はまだ明確には現れず、曲率は本稿の射程外である（§7・§8）。

\begin{center}\rule{0.5\linewidth}{0.5pt}\end{center}

\subsection{§1
規約（中／外の読み方）}\label{ux898fux7d04ux4e2dux5916ux306eux8aadux307fux65b9}

以下、形式の記述に入る。読解のため二つの規約を置く。

\begin{enumerate}
\def\labelenumi{\arabic{enumi}.}
\tightlist
\item
  【タグ】本稿の各言明には、それが系の中（不可知・無名の内部視点）か系の外（観測者の視点）かを原則として明示する。双対のもとで不変に保たれる錨は【不動点】【保存量】【離散構造】として示し、実在の所在は主張しない。
\item
  【動詞の規律】系の中の区別不能を述べるとき「〜が無い／存在しない」とは書かず、「系の中では〜と区別できない」と書く。「無い」と書いてよいのは離散構造・不動点・保存量についての客観的言明に限り、系の中の区別不能には用いない。
\end{enumerate}

【宣言】以下の公理および導出のうち、0・1・2
に関する言明は、すべて系の中の区別可能性についての言明であり、整数構造（離散・不動点・保存量）を否定しない。

\begin{center}\rule{0.5\linewidth}{0.5pt}\end{center}

\subsection{§2 公理（仮定）}\label{ux516cux7406ux4eeeux5b9a}

本系は以下の公理に立脚する。外部前提や過去文献への参照を用いない。物理的解釈は分離し、構造の言明に留める。

\textbf{公理1（差異の存在）：}
差異がある。「別な量がある」ということだけが言える。何が違うか・どれとどれが違うか・いくつあるかは言えない（無名）。差異を語ること自体が差異を前提とするため、これは自己言及的である。

\textbf{公理2（無名性の限界＝2）：}
【系の中】差異が立つ最小は二である。無名性ゆえ上限も二である。系の中では
0・1・2 で閉じ、2 が天井となる（§0.5 の導出による：0 と 1
は区別できず、2 と 2 以上も区別できない）。3 以上を 3
以上として扱うことは数え上げ＝ラベル＝次元の発生であり、無名性の破れである。これは系の中の区別可能性の限界であって、整数構造としての
3 以上を排除しない。

\textbf{公理3（揺らぎ）：} 各量は固有の揺らぎ δ を持つ。δ
はその量の属性でスカラーである。揺らぎの大きさ ±1/2
は導出ではなく措定である（§6）。

\textbf{公理4（積の保存：在る／無いの縁）：}
【系の中】二つの揺らぎの積は保存する。δ\textsubscript{1}δ\textsubscript{2}=1/2。これは「在る／無い」という存在の二値の、その縁の不変量（積・双曲
SO(1,1)・面積）であり、後述する複素内積の虚部 ω に対応する。
この虚部というのも、本来は無名性の原則から選択問題であり、どちらを虚部ととるかは任意である。

\textbf{公理5（二乗和の保存：差異の束ね）：}
【系の中】系は有界かつ対称（ラベルの読み替えに不変な無名性）である。これにより二者を対等に束ねる不変量が立つ。Σν²=R²（二乗和・SO(2)対称・スケール
R＝距離 g）。ここで SO(2)
は二乗和を保つ対称性を指すのみで、円・球面という幾何像を仮定するものではない。これは「別な量がある」という差異の存在の不変量であり、後述する複素内積の実部
g に対応する。
この実部というのも、本来は無名性の原則から選択問題であり、どちらを実部ととるかは任意である。

\textbf{補題（閉性は公理ではなく無名性の帰結）：}
【系の中】端（境界）は片側にしか隣を持たない特権的な点であり、無名性に反する。ゆえに無名性を保つ系は端を持てず、閉じる------しかも周期的（回転対称）に閉じる。よって「閉じた系」は独立な公理ではなく、無名性（公理1・2）の帰結として導かれる。旧版で独立公理として置いていた閉性の仮定は、本改訂版で除いた。

\textbf{注（公理4 と公理5 の相互導出不能）：}
【系の中】積（公理4・ω）と二乗和（公理5・g）は相互に導出できない。これは技術的事情ではなく、「在るか」と「違うか」が無名性の下で独立に立つ二つの根源区別だからである。一方から他方が出たら、無名性の限界が
2 より広いことになってしまう。

\begin{center}\rule{0.5\linewidth}{0.5pt}\end{center}

\subsection{§3
離散性の導出（離散性は公理ではなく導出量である）}\label{ux96e2ux6563ux6027ux306eux5c0eux51faux96e2ux6563ux6027ux306fux516cux7406ux3067ux306fux306aux304fux5c0eux51faux91cfux3067ux3042ux308b}

本系の離散性は公理として置かれていない。二段で導出される。

\textbf{第一段（根本的離散性：公理2 から）。} 【系の中】系の中では 1
が単独で立たない（公理2）。すなわち連続体を細分するための単位「1」が、対照を欠いて
0
と区別できない。連続濃度は「単独で立つ単位を無限に細分できる」ことで初めて成立するが、その単位が立たない。ゆえに系の中で構成できるのは二値の有無判定の反復＝二倍・二分（×2,
÷2）のみであり、計数は離散（整数）となる。底が 2
であること、整数性、後述の半ビットは、いずれもこの「1
が単独で立たない」から降りてくる。

\textbf{第二段（モード番号の離散性：有界性から）。} 【離散構造】公理5
の有界性は、§2
の補題により\textbf{周期的に閉じた系}（端を持たず一周して戻る系）として現れる。端のある弦
{[}0,L{]}
ではない点が、補題（端は無名性に反する）と整合する。周期的に閉じた系（周長
L）の上で定常波が立つには周期境界条件を満たさねばならず、許される波長は周長の整数分の一に量子化される（λ\_n=L/n、n=1,2,3,\ldots）。スペクトルが整数
n で番号づけられる。この整数モード番号 n
の発生が、数え上げ＝次元の生成であり、公理2 の天井（2）を破って 3
以上へ抜ける機構である。

すなわち、離散性は公理5 が公理1〜4
に対してすることの名であって、独立な原理ではない。

\begin{center}\rule{0.5\linewidth}{0.5pt}\end{center}

\subsection{§4 基本量 N
の導入}\label{ux57faux672cux91cf-n-ux306eux5c0eux5165}

【保存量】保存される量は、観測される一乗量 ν, λ ではなく二乗量
Σν²=R²（公理5）である。一乗量 ν, λ は共役関係 νλ=k
の下で互いに逆数的に振れる観測量であり、その積の背後に不変な二乗量が保存される。

基本量を、二乗量の対数として導入する。底は公理2（系の中の二値性）により
2 を採る。係数 1/2 を付して整数列とする：

\[N \equiv \tfrac{1}{2}\log_2(\nu^2) = \log_2\nu .\]

【不動点】ν=2\textsuperscript{n} のとき N=n（整数）。N=0 は
ν=1（自己双対の不動点）に対応する。【系の中】系の中では ν=1 は ν=0
と区別できないが（公理2）、これは不動点・整数構造としての 1
を否定しない。N は 0 を中心に正負対称に整数全体 $\mathbb{Z}$ を表現する。N
の符号は共役対のいずれの側から見るかに対応する。

【系の中／外】虚数単位 i
は、90度位相差を持つ二実数成分を複素記法で畳んだ記号に過ぎず、本系の公理は実数で閉じる（複素数は導入しない）。後述（§7
注）の通り、複素構造はこの実数二者性の帰結として現れる。

\begin{center}\rule{0.5\linewidth}{0.5pt}\end{center}

\subsection{§5 共役関係の加法的恒等式（g
側）}\label{ux5171ux5f79ux95a2ux4fc2ux306eux52a0ux6cd5ux7684ux6052ux7b49ux5f0fg-ux5074}

【系の外（観測者）】共役関係 νλ=k に底 2 の対数を取ると、積が和になる：

\[\nu\lambda=k \;\xrightarrow{\ \log_2\ }\; N_\nu + N_\lambda = K \quad(K\equiv\log_2 k).\]

これは恒等式であり、保存量は K（共役対の N の和）である。三項 N\_ν,
N\_λ, K のいずれを右辺に選んでも

\[N_\lambda = K - N_\nu,\qquad N_\nu = K - N_\lambda,\qquad N_\nu + N_\lambda = K\]

と恒等式のまま保たれ、入れ替えで生じるのは対数の引き算（除算）というラベルの読み替えのみである。【系の中】したがって、いずれの量を特権化することもなく、無名性が維持される。

【系の外（観測者）】スケールの変更（k の変更）は N
軸の平行移動に対応し、スケール不変性が並進対称性として現れる。k=1 では
K=0 となり N\_λ=−N\_ν（0 を中心とする符号反転対称）である。

\textbf{注（恒等式は構造を生まない）：}
この加法恒等式は対数の定義そのものであり、それ自体は新しい構造を生まない。構造を生むのは、次節の揺らぎ
Δ（ω 側）である。本恒等式は g 側（距離・二乗和）の言い換えにすぎない。

\begin{center}\rule{0.5\linewidth}{0.5pt}\end{center}

\subsection{§6 揺らぎ Δ と不確定性（ω
側、構造を生む側）}\label{ux63faux3089ux304e-ux3b4-ux3068ux4e0dux78baux5b9aux6027ux3c9-ux5074ux69cbux9020ux3092ux751fux3080ux5074}

構造を生むのは揺らぎ Δ
である。ここを正確に置くため、まず誤読を二つ殺す。

\subsubsection{6.1
隠れ変数・観測誤差の否定}\label{ux96a0ux308cux5909ux6570ux89b3ux6e2cux8aa4ux5deeux306eux5426ux5b9a}

【離散構造】ν は離散（整数値）であり、連続な小数部を持たない。「本当は
ν=2.0±δ
の連続値があり、観測精度の限界で整数に見える」のではなく、系が離散だから整数である。これは実在の主張ではなく、離散構造の帰結である（§0.5
の注）。

ν=2 とは、ν が整数のまま値を揺らがせ、その中心値（中央値）が 2
であることを意味する。たとえば値は近傍の整数 1, 2, 3 を取り得て、中央が
2。値そのものが整数の間を飛ぶ。これは系に内在する揺らぎであって、固定した真値に乗る観測誤差ではない。

【系の中】系の中の量自身は整数である。ゆえに 0.5
のような半端な値を取り得ない。系の中で揺らぎが現れるとすれば、整数きざみ、すなわち
±1 のずれとしてのみ現れる。

\subsubsection{6.2 ±1/2
は導出ではなく措定である}\label{ux306fux5c0eux51faux3067ux306fux306aux304fux63aaux5b9aux3067ux3042ux308b}

【否定】ここで「標準偏差は
±1/2」と言いたくなる。だが本稿はこれを導出しない。そもそもこの系には、標準偏差という概念を組み立てる前提が無い。

【理由】標準偏差を定義するには、同一の対象を繰り返し観測し、そのばらつきを集計する操作が要る。だが------

\begin{itemize}
\tightlist
\item
  【系の中】無名性により、観測を互いに個別化（これは一回目、これは二回目）できない。
\item
  時間発展を定義していないので、観測を並べる順序が無い。
\item
  観測の間で真の値が同一に保たれるとも仮定していない。
\end{itemize}

ゆえに「複数回観測する」という言葉自体が、この系では自己矛盾である。その一言が、個別化・時間順序・値の持続という隠れた前提を多数持ち込んでしまう。

【措定】したがって ±1/2
は、計算で得た値ではなく措定（こうであろうという仮定）である。どの計算方法で
1/2
を導いても、必ずいずれかの隠れた前提（個別化・時間順序・持続）を持ち込むことになる。標準偏差を厳密に計算できないのではなく、標準偏差という概念を組み立てる前提自体が無い。本稿は
±1/2 を、導出ではなく仮定として置く。

\subsubsection{6.3
揺らぎの保存則（和ゼロが定義、積一定が導出）}\label{ux63faux3089ux304eux306eux4fddux5b58ux5247ux548cux30bcux30edux304cux5b9aux7fa9ux7a4dux4e00ux5b9aux304cux5c0eux51fa}

\textbf{注（§6.1 と §6.3 の揺らぎは別物である）：} §6.1
で述べた「系の中の整数量が整数きざみ（±1）で揺らぐ」ことと、本節 §6.3 で
N 空間に置く半ビット揺らぎ Δ=±1/2
は、\textbf{別個の措定であり、一方から他方を導くものではない}。両者の間に関連（同一性・導出・近似のいずれも）があるとは主張しない。§6.1
は系の中の整数量の揺らぎについての言明、§6.3 は基本量 N
の側に措定する揺らぎについての言明であり、本稿は両者を独立した仮定として扱い、その整合や統一は今後の課題とする。繋がらないものを繋げば必ず隠れた前提が入る（§6.2
の ±1/2 と同じ作法）ため、本稿は繋げない。

【系の中】基本量 N の側に揺らぎを置く（N
空間で揺らぎを措定すると、観測量空間では揺らぎがスケール依存となり、整数計数の揺らぎとして整合的である）。各基本量に揺らぎ
Δ=±1/2（半ビット、措定、スケール不変）を与える。共役対は和の保存
N\_a+N\_b=K を満たし、揺らぎは

\[\boxed{\ \Delta_a + \Delta_b = 0\ }\]

を満たす（これを定義とする）。

【系の外（観測者）】逆写像 a=2\^{}\{N\_a+Δ\_a\} により、

\[k = ab = 2^{(N_a+N_b)+(\Delta_a+\Delta_b)} = 2^{K+0} = 2^K \quad(\text{一定})\]

が導かれ、観測量の揺らぎの積は一定（最小不確定の下限）となる。すなわち、観測量空間における不確定性関係（積一定）は、N
空間における揺らぎの和ゼロから導出される。不確定性の最小単位 1/2
は、措定された半ビット ±1/2 に対応する。なお単一のゆらぎ ±1/2
と、二つのゆらぎの積 δ\textsubscript{1}δ\textsubscript{2}=1/2（公理4）は同一の量ではなく、同じ
ω（面積）側に属する量として対応づけられる。

\begin{center}\rule{0.5\linewidth}{0.5pt}\end{center}

\subsection{§7 主結果：スケールの対数写像と g/ω
分離（限定された極限）}\label{ux4e3bux7d50ux679cux30b9ux30b1ux30fcux30ebux306eux5bfeux6570ux5199ux50cfux3068-gux3c9-ux5206ux96e2ux9650ux5b9aux3055ux308cux305fux6975ux9650}

\textbf{【注：ここでの「スケール」について】}
本節で平坦化の対象とするのは Σν²=R² の
R、すなわち系のスケール（大きさ）である。これは曲率テンソルでも球面の半径でもない。Σν²=R²
は二乗和が R²
に等しいという保存則であって、円や球面の幾何を要求しない。本稿は「曲率の平坦化」を主張するのではなく「スケールの対数による線形化（平坦化）」を述べる。曲率テンソルへの拡張は本稿の射程外である（§8・§9）。本稿は、この系の連続極限が曲率テンソルへ拡張できることを主張せず、むしろ離散系においてのみ成立する可能性が高いと見ている。

\subsubsection{7.1
二空間における線形化の相補性}\label{ux4e8cux7a7aux9593ux306bux304aux3051ux308bux7ddaux5f62ux5316ux306eux76f8ux88dcux6027}

【系の外（観測者）】公理5 の保存量（系の中の不変量）は、観測者には二乗和
Σν²=R²（g・スケール
R）として現れ、これがスケールを表す保存則となる。共役関係 νλ=k
は積（非線形）である。対数による N
空間においては、この関係が反転する。νλ=k
は加法的（N\_ν+N\_λ=K、§5）すなわち平坦な超平面となるが、二乗和 Σν²
は一般には指数の和 Σ2\^{}\{2N\_n\} となり、N について線形化されない。

\subsubsection{7.2
高対称・無名性の極限での平坦化}\label{ux9ad8ux5bfeux79f0ux7121ux540dux6027ux306eux6975ux9650ux3067ux306eux5e73ux5766ux5316}

【系の中】全対称（どの軸も右辺に来れる、すなわちすべての ν\_n²
が等しい）の極限を考える。この極限は、すべての ν\_n² が等しい値 v
を取ることに帰着する（軸数は実質 2 に縮約される。これは公理2
の二値性の再現でもある）。このとき Σν²=（軸数）·v
となり、対数が和に作用可能となる：

\[\log_2\!\left(\sum\nu^2\right) = \log_2(\text{軸数}) + \log_2 v.\]

すなわち、スケールを表す二乗和が、平坦な一乗の加法関係に写像される：

\[\sum\nu^2 = R^2 \;\xrightarrow{\ \log_2,\ \text{全対称}\ }\; N + \tfrac12\log_2(\text{軸数}) = N_R \quad(\text{恒等式、平坦}).\]

スケールの源は二乗（冪）であり、対数は冪を線形化する。さらに本系はスケール不変であるため、スケール依存性が並進に沿った直線として平坦化される。この極限において、§9
で述べる Weyl
変換が残す微分項に対応する非一様項は、全対称（空間的一様）の仮定により消失する。

ただし全対称極限は、平坦化を厳密に成立させる代償として、g
側の自由度をほぼ消尽する。すなわち §7.2
の平坦化そのものの内実は薄い。本稿の主張の実質は、平坦化そのものより、次節
§7.3 の g/ω 分離------g が一値に縮退してもなお
ω（量子化）が残ること------にある。

\subsubsection{7.3
写像の核心：スケール（g）と量子化（ω）の分離}\label{ux5199ux50cfux306eux6838ux5fc3ux30b9ux30b1ux30fcux30ebgux3068ux91cfux5b50ux5316ux3c9ux306eux5206ux96e2}

ここが本稿の核心である。上の平坦化は「スケールの情報を捨てている」のではない。

【系の外（観測者）】公理5 の保存量は、観測者には二乗和
Σν²、すなわち距離・計量 g（後述の複素内積の実部）として現れ、Born
ノルム・規格化の側である。対数写像はこの g
側に作用し、全対称極限でそれを平坦な座標へ畳む。g
側はスケールを担うが、量子化を担わない。スケールは g
に丸ごと宿り、平坦化可能（gauge として除ける）。

【系の中】一方、揺らぎ
Δ＝δ\textsubscript{1}δ\textsubscript{2}=1/2（公理4）は面積・シンプレクティック形式
ω（複素内積の虚部）であり、量子化を担う側である。この ω
は対数写像で潰れず、平坦な N 面上に半ビット ±1/2
の不変量として残る。全対称極限で g（平均・スケール）が一値 v
に縮退しても、ω（揺らぎ Δ）は縮退しない。

すなわち対数写像は、\textbf{スケール（g・gauge）と量子化（ω・面積・不変量）を分離する写像}である。ここで「分離」とは
operational に、写像が g（スケール・二乗和）を平坦化する一方、ω（積
δ\textsubscript{1}δ\textsubscript{2}）を不変に保つことを指す。これはシンプレクティック形式を厳密に保つ正準変換であるという主張ではない（その厳密化は今後の課題）。スケールを平坦面に押し込んで
gauge
化し、量子化をその平坦面上の不変量として純化する。これが「全対称極限で平坦化しても物理が消えない」ことの内実である------物理（量子化）は初めから
g（平均）には無く、ω（揺らぎ Δ）に宿っていた。

\subsubsection{7.4 写像の構造}\label{ux5199ux50cfux306eux69cbux9020}

【系の外（観測者）】自己双対の不動点 k=1
かつ全対称・無名性の極限において：

\begin{itemize}
\tightlist
\item
  観測量空間（ν 空間）のスケールを持つ系（二乗和 Σν²、スケール R、g
  側）は、
\item
  対数接続により N 空間の平坦な超平面（一乗和 ΣN=K
  のアフィン構造）に写像され、
\item
  逆写像 ν=2\^{}N により観測量空間に厳密に復元される。
\item
  そして揺らぎ Δ（ω 側）は、この写像の前後で不変に保たれる。
\end{itemize}

\textbf{注（複素構造は実数二者性の帰結）：} 【系の中／外】ν\textsubscript{1}, ν\textsubscript{2} を動径
ν と位相 θ に読み替えると ν\textsubscript{1}=ν cosθ, ν\textsubscript{2}=ν sinθ となり、二者は z=ν
e\^{}\{iθ\} の極座標表示と幾何的に同型である。二乗和
ν²=ν\textsubscript{1}²+ν\textsubscript{2}²=\textbar z\textbar² が実部 g、揺らぎの面積 δ\textsubscript{1}δ\textsubscript{2}=1/2 が虚部
ω、cos↔sin の回転が複素構造 J
に対応する。複素数は公理に置かれておらず、実数二者性（公理2
の「2」）と二乗和（公理5）の帰結として現れる。

\begin{center}\rule{0.5\linewidth}{0.5pt}\end{center}

\subsection{§8
限定事項（Limitations）}\label{ux9650ux5b9aux4e8bux9805limitations}

本結果の適用範囲は厳格に限定される。

\begin{enumerate}
\def\labelenumi{\arabic{enumi}.}
\item
  \textbf{本結果は特殊解であり、一般解ではない。} 成立は
  k=1（自己双対の不動点）かつ全対称・無名性の極限に限られる。この極限を離れる（k≠1、非一様、非対称）と、§7.2
  で消失した非一様項（Weyl
  変換の微分項に対応）が一般には復活し、平坦化は成立しない。
\item
  \textbf{±1/2 は措定であり、導出されていない。} §6.2
  の通り、この系には標準偏差という概念を組み立てる前提が無い。±1/2
  を導く全ての計算は隠れた前提を持ち込む。なお本稿の主結果（§7.3 の g/ω
  分離）は、この ±1/2 の措定を前提とする。措定が外れれば結論も変わる。
\item
  \textbf{反証の網羅的検証は行えていない。}
  本結果が反例を持たないことは示されていない。
\item
  \textbf{既知枠組みとの関係は未確定。} 本結果が Weyl
  変換・共形平坦性・共形宇宙論等の既知枠組みの特殊例に留まるのか、それを超えるのかは現時点で確定できない。特に、§7.3
  の g/ω 分離が既知の Kähler
  幾何・共形平坦性とどう異同するかは、厳密には対照されていない。
\item
  \textbf{本稿はスケール R のみを扱い、曲率テンソルは扱わない。} Σν²=R²
  の R
  は系のスケール（大きさ）であって、曲率テンソルでも球面の半径でもない。微分構造（曲率テンソル）まで含めた厳密証明は与えておらず、曲率・一般の歪んだ時空への拡張は本稿の射程外である。
\item
  \textbf{物理的同定を行わない。}
  本系は代数的・幾何的（情報的）構造であり、観測量・基本量への物理的解釈（時間・空間・エネルギー・粒子等）は付さない。物理理論との同型は別途の課題である。
\item
  \textbf{二つの揺らぎの関連・統一を与えていない。} §6.1
  の整数きざみ揺らぎ（±1）と §6.3 の半ビット揺らぎ Δ=±1/2
  は別個の措定であり、両者の関連（同一性・導出・近似）は主張しない。N
  空間に揺らぎを置くことの、系の中の整数揺らぎからの導出は本稿では行わない。統一は今後の課題とする。
\end{enumerate}

\begin{center}\rule{0.5\linewidth}{0.5pt}\end{center}

\subsection{§9
既存枠組みとの位置づけ（参考）}\label{ux65e2ux5b58ux67a0ux7d44ux307fux3068ux306eux4f4dux7f6eux3065ux3051ux53c2ux8003}

曲がった時空における物理量の計算は一般に困難である。既存のアプローチの一つに
Weyl 変換 \(g_{\mu\nu}\to\Omega^2 g_{\mu\nu}\) があるが、標準的な作用は
Weyl 再スケールの下で不変ではなく、曲率スカラーは

\[R \to e^{-2\sigma}\left(R - 6\,\Box\sigma - 6\,g^{\mu\nu}\partial_\mu\sigma\,\partial_\nu\sigma\right)\]

と変換し、曲率の情報は再スケール関数 σ
の微分項として残る。すなわち平坦化しても曲率は完全には消えない。また共形平坦性は二次元では局所的に常に成立するが、高次元の一般の計量では成立しない。

本報告は、これらの既存枠組みとは異なる出発点------ごく少数の公理のみ------から、ある限定された極限において、スケールを表す量（曲率テンソルではない）が微分項を残さず平坦な構造に写像され、かつ量子化（ω）がその写像で不変に保たれる、という構造を提示するものである。本稿が扱うのは曲率テンソルを再スケールする
Weyl 変換とは射程が異なり、スケールの対数化である。Weyl
の焼き直しでもなければ、曲率を扱うという主張でもない。曲率への拡張は今後の課題とする。一般化は主張しない。既知枠組みの特殊例である可能性も排除しない。

\begin{center}\rule{0.5\linewidth}{0.5pt}\end{center}

\subsection{§10 結語}\label{ux7d50ux8a9e}

本報告は、ごく少数の公理（差異・無名性の限界＝2・揺らぎ・積の保存・二乗和の保存）のみから出発し、系の中（不可知）と系の外（観測）を区別し、客観的な錨を実在ではなく不動点・保存量に置いた上で、基本量を二乗量の対数
N=log\textsubscript{2}ν
として導入したとき、極めて高い対称性と無名性の極限において、スケールを表す二乗和（Σν²=R²・g）が対数接続により平坦な超平面（アフィン構造）に写像され得ること、そしてその写像がスケール（g）と量子化（ω・揺らぎ
Δ）を分離し、量子化を平坦面上の不変量として保つこと、を限定された構造として記録するものである。ここで
R は系のスケールであって曲率テンソルではない。

これを一般的事実の発見として主張せず、限定範囲での検証結果として、また
±1/2
の措定性と反証検証の未完を明記した上で、記録する。本構造が既知枠組みの特殊例であるか否か、一般化が可能であるか否かは、今後の検討に委ねられる。

\begin{center}\rule{0.5\linewidth}{0.5pt}\end{center}

\emph{本報告は予備的・限定的記録であり、結果の一般性・厳密性・新規性のいずれについても確定的主張を含まない。}
\end{document}
